\chapter{Background and Related Work}
\label{ch:bg}

\chapterhandoff{Chapter~\ref{ch:intro} argued that choosing how to show the
page, rather than fixing that choice once, is worth studying if three
conditions hold. This chapter establishes what is already known about how web
agents see pages, about spending computation selectively in other settings,
and about how efficiency is accounted for. It then identifies the one decision
that none of that existing work is trying to make well.}

\section{How a web agent acts}
\label{sec:agent-loop}

A web agent operates a browser through a loop: it is shown the page, picks one
action, the browser performs it, and it is shown the page again. Written out,
at step $t$ the agent receives an \emph{observation} $o_t$ describing the
current page, selects an \emph{action} $a_t$ from a fixed list of permitted
browser operations (clicking, typing, selecting, scrolling, navigating,
waiting, managing tabs, and declaring the task finished), executes it against
the browser, and receives
the next observation. The loop terminates when the agent emits a completion
action, exhausts a step budget, or errors out. A benchmark \emph{evaluator} then
inspects the final environment state and returns a binary \emph{task success}.

The loop has two properties that shape everything which follows. First, the
observation is \emph{constructed}: something has to decide which parts of the
browser's internal state become text, which become pixels, and how they are
framed for the model. Second, the loop is \emph{long} (tens of steps for a
typical task), so the benefit of any per-step decision is spread across the
whole sequence, while its cost is multiplied by the length of that sequence.

\section{Ways to represent a web page}
\label{sec:representations}

Our six modes are assembled from the same building blocks that existing
systems use to describe a page to a model.

\paragraph{Structured text.} The browser's accessibility tree (AXTree) exposes
elements, roles, names, states and identifiers. It is a compact, textual, and
inherently \emph{selective} view: it names what is interactable and discards
layout, colour and imagery. Systems that consume it must decide how much of it
to include, which is itself an active research question
\citep{deng2023mind2web, lu2024weblinx}.

\paragraph{Screenshots.} A rendered image preserves everything the tree
discards, at the cost of requiring the model to localise elements visually. The
difficulty of turning a visual intention into a valid browser action (\emph{grounding}) is well documented; \citet{zheng2024gpt4v}
show that a capable model able to read both text and images becomes a
competent web agent largely once that grounding problem is solved for it.

\paragraph{Set-of-Marks.} \citet{yang2023som} overlay numbered marks on
candidate elements so the model can refer to a region by index rather than by
coordinate. In web-agent systems this image is typically accompanied by a text
listing of the same numbered elements, so a ``Set-of-Marks observation'' is in
practice three things at once: an annotated image, a listing of the marks, and
a prompt explaining how the numbering works. Chapter~\ref{ch:system} takes that bundle apart. That
it is a bundle of three separable things, rather than one setting, is what
makes the comparison in this dissertation possible at all.

\subsection{Why the choice changes both capability and cost}

What a representation is worth and what it costs move independently of each
other. Adding an image can supply information no text encoding carries, such as visual attributes, spatial relationships, or rendered content absent from the tree, and can equally supply nothing at all if the
task's discriminating information was already in the text. Meanwhile the cost
is paid regardless: the image is charged for as input, encoding it takes time,
and a longer input makes the model slower to respond. A representation is therefore worth
its price only conditionally, and the condition is a property of the task and
the interaction state rather than of the representation.

\section{Selective computation: four programmes, four decision variables}
\label{sec:selective}

Selective computation is a mature idea with several literatures. Organising them
by \emph{what each one chooses} rather than chronologically makes the gap this dissertation addresses visible as a structural property rather than as
an assertion (Figure~\ref{fig:f2}).

\begin{figure}[t]
  \centering
  \includegraphics[width=\textwidth]{fig_f2_literature_map.pdf}
  \caption[Literature map, organised by decision variable]{%
  \textbf{Twenty-seven works, clustered by the variable each one selects.} Not a
  timeline: the organising axis is \emph{what the method chooses}, because that
  is the axis on which this dissertation differs from all four clusters.
  Observation-representation work (left) selects which elements enter a fixed
  encoding; model routing (top) selects which model, leaving the input encoding
  untouched; confidence-based deferral (right) selects whether to answer, using
  a signal that typically requires paying for inference first; adaptive
  inference (bottom) selects how much computation to spend, which is depth or
  width rather than what is looked at. The centre is empty: for a
  \emph{fixed} agent, choosing \emph{which page encoding} to buy at a step,
  using signals available \emph{before} paying for it. Cluster summaries are our
  synthesis, not quotations.}
  \label{fig:f2}
\end{figure}

\paragraph{Cluster 1: which elements enter the observation.}
\citet{deng2023mind2web} rank and prune HTML elements before action prediction;
\citet{lu2024weblinx} learn a component that ranks page elements by likely
relevance and keeps only the top ones. Both
reduce cost inside one representation. The alternative they never consider is a
\emph{different} representation. Benchmarks \citep{zhou2024webarena,
koh2024visualwebarena} and end-to-end agents \citep{he2024webvoyager,
zheng2024gpt4v} fix the representation as a configuration and optimise the
policy that consumes it.

\paragraph{Cluster 2: which model answers.} FrugalGPT \citep{chen2024frugalgpt},
AutoMix \citep{aggarwal2024automix}, Hybrid LLM \citep{ding2024hybridllm},
RouteLLM \citep{ong2025routellm}, RouterDC \citep{chen2024routerdc} and
TensorOpera Router \citep{stripelis2024tensoropera} choose between models, or
try them in order from cheapest to most capable. Structurally this is the closest analogue to our problem:
a cheap decision made before an expensive call, trained on outcome labels.
The switched object, however, is model identity, and the input the model
receives is held constant. Several of these methods additionally gate on a
signal produced by \emph{running} the cheap model first
\citep{chen2024frugalgpt, aggarwal2024automix}, which our setting forbids by
construction.

\paragraph{Cluster 3: whether to answer at all.} Selective classification
\citep{geifman2017selective, geifman2019selectivenet, liu2019deepgamblers},
post-hoc confidence repair \citep{cattelan2024confidence} and learning-to-defer
\citep{mozannar2020defer} choose between predicting and abstaining or deferring.
The escalation endpoint is abstention or an external expert, not a richer view
of the same input, and the signal used to decide is normally read off a model
run that has already happened --- which is again downstream of the cost we are
trying to avoid paying.

\paragraph{Cluster 4: how much computation to spend.} Adaptive computation
time \citep{graves2016act}, early exit \citep{bolukbasi2017adaptive,
xin2020deebert, schuster2022calm}, token sparsification \citep{rao2021dynamicvit},
resolution cascades \citep{wang2021notallimages} and reasoning-budget prediction
\citep{han2025tokenbudget} allocate internal computation per input. The nearest
of these to our question is \citet{wang2021notallimages}, which escalates image resolution until confident. But
the signal that triggers the upgrade is the model's own confidence after the
cheap pass has already run, and the task is a single classification rather
than a sequence of browser steps.
\citet{han2025tokenbudget} is the nearest on the other axis: it predicts a
resource requirement \emph{ex ante} from cheap features, but the resource is
output-side reasoning length rather than input-side perception.

\subsection{Two neighbours that must be addressed directly}
\label{sec:neighbours}

Some recent work sits close enough that a general novelty claim would be overstated, and we state the distinction rather than the distance.

\citet{enomoto2026readmore} vary observation representation and thinking budget
together, and report how the best observation configuration shifts with model
capability and reasoning budget. This forbids the claim that prior web-agent
work has not studied adaptive observations. The difference is the granularity of
the decision: their result is a \emph{configuration-level} regularity (\emph{given this model and this budget, prefer this observation}), whereas the object here is a per-task, per-state decision for a \emph{fixed} agent, gated on
signals available before the richer observation is constructed. Their finding is
in fact a precondition of ours being interesting: if the best representation
never varied, there would be nothing to route.

\citet{montreuil2026defer} generalise learning-to-defer so that the action
selects both an expert and the additional information supplied to it, so that
the price of acquiring that extra information is part of the decision itself. That is, formally, the closest
statement of our decision problem that we are aware of. It is a one-shot
supervised formulation: it does not model sequential browser state, does not fix
the agent while varying the acquisition, and does not restrict the gate to
signals obtainable before acquisition. Our contribution relative to it is empirical rather than formal: we instantiate that decision in a real
sequential environment and measure what the terms are actually worth.

\section{Efficiency accounting}
\label{sec:bg-efficiency}

Cost claims in this area are frequently under-specified, and the correction is
well established. \citet{fernandez2025energy} show that inference energy
estimated from FLOPs or theoretical utilisation can diverge substantially from
measured energy, and that the divergence depends on workload, software stack and accelerator, so there is no hardware-independent token-to-energy conversion
to appeal to. \citet{lannelongue2021green} make the same point for GPU-hours,
which quantify activity rather than energy. \citet{strubell2019energy} is the
field's founding precedent for reporting these quantities at all, and
\citet{henderson2020systematic} the standard argument that they should be
reported systematically rather than anecdotally. \citet{luccioni2022bloom} and
\citet{gupta2022chasing} show that the reported number moves substantially
depending on which stages of energy use are counted --- manufacturing,
cooling, the data centre around the chip --- which is why
Section~\ref{sec:sustainability} names exactly what ours includes and declines
to call it a footprint.

A second accounting point is specific to routing and is easy to omit: the router
is part of the system. Feature extraction, the routing model's own inference,
any additional parsing or image preparation, and any retries caused by a wrong
routing decision are all system cost. A comparison that prices the routed
execution but not the routing is not a comparison of policies.

\section{The gap}
\label{sec:gap}

Existing web-agent work establishes that page representation materially affects
grounding, context cost and task performance, and recent work adds that the best
representation shifts with model capability and reasoning budget. Separately,
routing, selective prediction and adaptive-inference literatures establish that
computation can be allocated per input, using either a cheap ex-ante predictor
or a confidence signal computed downstream.

Neither line answers the question in the middle. For a \emph{fixed} agent at a
\emph{given} interaction state, can we predict, from information available
\emph{before} the expensive representation is paid for, how much extra a
richer page representation is worth over a cheaper one? The four clusters
supply, respectively, evidence that representations differ, machinery for
choosing between models, machinery for choosing whether to act, and machinery
for choosing how hard to think. What is missing is a measurement of whether
representation itself is a routable variable, and at what price.

That is an empirical question with three sequential parts, and it is answerable
only inside an experimental design that varies representation while holding
everything else constant.

\chapterhandoff{Prior work establishes that page representation changes both
capability and cost, and that selective computation is a viable design pattern
in adjacent settings. It does not supply a comparison in which the
representation is the only thing that changes, nor does it price what acting
on such a comparison would itself cost. Chapter~\ref{ch:system} therefore defines the
experimental system: what is varied, what is held fixed, and what is measured,
in enough detail that a hostile reader can reconstruct which differences are
attributable to representation and which are not.}
